\newpage
\section*{Resumen}
\begin{spacing}{1.4}

El desarrollo de sistemas de diagnóstico asistido por computadora mediante Inteligencia Artificial (IA) enfrenta un desafío crítico en cardiología: la escasez de datos etiquetados de alta calidad y el desbalance de clases en patologías raras. Esta investigación propone el desarrollo de un generador paramétrico de señales electrocardiográficas (ECG) sintéticas basado en la superposición de funciones Gaussianas asimétricas. A diferencia de los modelos de aprendizaje profundo (como las GANs), que carecen de control fisiológico y sufren de alucinaciones, este enfoque permite una parametrización explícita de la morfología P-QRS-T.

Mediante algoritmos de optimización numérica (L-BFGS-B) y utilizando registros de la base de datos MIT-BIH como semilla, el modelo calibra amplitudes, posiciones temporales y el nivel de la línea base para reproducir la morfología de señales reales. La validación se realizó sobre 2.237 latidos sinusales normales del registro 100 de dicha base, extraídos en una ventana asimétrica de 750 ms que contiene el complejo P-QRS-T completo, obteniendo una correlación de Pearson media de 0,9969 (desviación estándar 0,0015) y un RMSE medio de 0,0779 (desviación estándar 0,0148) en unidades de señal normalizada, con un sesgo del residuo indistinguible de cero.

Esta versión del trabajo incorpora además una segunda dimensión, ausente en la anterior: la organización temporal de la secuencia de latidos. Al modelo morfológico se le acopla un módulo de ritmo que produce la serie de intervalos entre latidos por síntesis espectral, con su exponente de escalamiento como parámetro de entrada declarado. El acoplamiento es selectivo en el tiempo: el complejo QRS no se estira con la frecuencia cardíaca, mientras que las ondas P y T sí. La validación de ese módulo no descansa en la indistinguibilidad respecto de una serie real, que fue el criterio del estado del arte, sino en la medición de una propiedad declarada: sobre una cohorte de cincuenta sujetos sintéticos de veinticuatro horas, el exponente pedido se recupera con un error absoluto medio de 0,014, y el sesgo del estimador queda caracterizado en función del largo del registro.

Estas cifras corresponden a la segunda iteración del modelo. Una primera implementación de veinte parámetros alcanzó una correlación de 0,9647 y un RMSE de 0,3985, y el análisis de su error residual identificó dos defectos concretos: la saturación del límite superior de amplitud de la onda R y la ausencia de un término de línea base. La corrección de ambos redujo el RMSE en un 87,6\% y llevó el sesgo sistemático del residuo a cero. El diagnóstico del residuo, y no solo la métrica agregada, resultó ser el instrumento que permitió esa mejora.

Los resultados obtenidos confirman la viabilidad del enfoque paramétrico para reproducir la morfología del ciclo P-QRS-T en ritmo sinusal normal. La incorporación de un motor de reglas basado en estándares clínicos (Rijnbeek/AHA) para simular variabilidad demográfica y patologías específicas se plantea como la extensión inmediata del trabajo, y no forma parte del alcance validado en esta etapa.

\end{spacing}
